
\documentclass{article} % For LaTeX2e
\usepackage{iclr2026_conference,times}

% Optional math commands from https://github.com/goodfeli/dlbook_notation.
\input{math_commands.tex}

\usepackage{hyperref}
\usepackage{url}
\usepackage{graphicx}
\usepackage{booktabs} 

\title{The Noise Accumulation Bottleneck: Retrieval Mechanisms Outperform Long Context Scaling in Time Series Forecasting}

% Authors must not appear in the submitted version. They should be hidden
% as long as the \iclrfinalcopy macro remains commented out below.
% Non-anonymous submissions will be rejected without review.

\author{Antiquus S.~Hippocampus, Natalia Cerebro \& Amelie P. Amygdale \thanks{ Use footnote for providing further information
about author (webpage, alternative address)---\emph{not} for acknowledging
funding agencies.  Funding acknowledgements go at the end of the paper.} \\
Department of Computer Science\\
Cranberry-Lemon University\\
Pittsburgh, PA 15213, USA \\
\texttt{\{hippo,brain,jen\}@cs.cranberry-lemon.edu} \\
\And
Ji Q. Ren \& Yevgeny LeNet \\
Department of Computational Neuroscience \\
University of the Witwatersrand \\
Joburg, South Africa \\
\texttt{\{robot,net\}@wits.ac.za} \\
\AND
Coauthor \\
Affiliation \\
Address \\
\texttt{email}
}

% The \author macro works with any number of authors. There are two commands
% used to separate the names and addresses of multiple authors: \And and \AND.
%
% Using \And between authors leaves it to \LaTeX{} to determine where to break
% the lines. Using \AND forces a linebreak at that point. So, if \LaTeX{}
% puts 3 of 4 authors names on the first line, and the last on the second
% line, try using \AND instead of \And before the third author name.

\newcommand{\fix}{\marginpar{FIX}}
\newcommand{\new}{\marginpar{NEW}}

%\iclrfinalcopy % Uncomment for camera-ready version, but NOT for submission.
\begin{document}


\maketitle

\begin{abstract}
Time Series Foundation Models (TSFMs) are increasingly adopting the long context paradigm from Natural Language Processing, assuming that long historical windows consistently improve predictive reasoning. However, this scaling hypothesis has not been tested in stochastic domains, where distant history often contains high-frequency noise rather than semantic signals. The proposed work challenges this assumption by examining the continuous context architectures, including PatchTST, on the ETTh1 benchmark. The analysis reveals a distinct inverse scaling law: the forecasting error increases as the length of the context increases. In particular, the performance degrades by over 68\% for a 3,000-step window, suggesting that the attention mechanisms do not filter out irrelevant historical volatility. The study assesses Retrieval-Augmented Forecasting (RAFT) as an alternative. Our results demonstrate that RAFT, using a fixed 720-step window with selective retrieval, achieves a mean squared error (MSE) of 0.379, outperforming the computationally costly long-context models (MSE 0.647). Moreover, the retrieval mechanism functions as a context-informed inductive bias, injecting essential historical signals as dynamic exogenous variables. Therefore, selective retrieval constitutes a necessary architectural shift for future foundation models.
\end{abstract}

\section{Introduction}

The recent surge in Time Series Foundation Models (TSFMs) has been fueled by "scaling hypothesis" - the assumption that larger models and extended context windows monotonically improve predictive reasoning. Architectures like Moirai \citep{woo2024unified} and Chronos \citep{ansari2024chronos} have shown that scaling effectively opens up zero-shot capabilities across variety of areas. But the critical question about the inductive bias of these architectures remains: \textit{Does the "Long Context" reasoning that drives Large Language Models (LLMs) translate directly to the stochastic domain of time series data?}

While natural language processing benefits from massive context because distant tokens often retain a high degree of semantic fidelity, time series data shows significant differences. Historical numerical sequences are often dominated by stochastic volatility rather than by semantic signals. Hypothetically, indiscriminate scaling of the context length in forecasting leads to Stochastic Noise Accumulation, where the focus of the attention mechanism is diluted by irrelevant history, degrading performance rather than improving it.

The literature is increasingly emphasizing the necessity for auxiliary information rather than a simple window extension. \cite{williams2024context} demonstrates that "context is key," showing that the forecasting accuracy increases significantly when models are given key details on basic textual descriptions. Similarly, \cite{arango2025chronosx} highlights the critical role of integrating exogenous variables to adapt pretrained models to subsequent tasks. Building on these findings, the proposed work argues that \textit{selective retrieval}, rather than a brute force window extension, is the optimal mechanism for providing this essential context.

The proposed work empirically tests this hypothesis on the ETTh1 benchmark, comparing vanilla Transformers, the state-of-the-art PatchTST \citep{nie2023time}, and a Retrieval-Augmented Forecasting (RAFT) baseline. The experimental results reveals a Inverse Scaling Law for stochastic time series:
\begin{itemize}
    \item \textbf{Long Context Fails:} Increasing the context from 720 to 3,000 tokens decreases PatchTST performance by 68\%, confirming that even state-of-the-art patching mechanisms can not filter out excessive historical noise.
    \item \textbf{Retrieval Wins:} A RAFT mechanism \citep{han2025raft} using only a 720-step window with selective retrieval achieves best performance (MSE 0.379), highly outperforming the computationally expensive 3,000-step models (MSE 0.647).
\end{itemize}

The analysis suggests that retrieved historical segments function as \textit{dynamic exogenous variables}, providing the context-informed inductive bias required to stabilize the forecasts in high-variance domains.

\section{Related Work}

\subsection{Time Series Foundation Models and Context Scaling}
Time Series Foundation Models (TSFMs) such as Moirai and Chronos leverage massive pretraining on diverse corpora to enable remarkable zero-shot generalization across distinct domains. By learning universal temporal patterns, these architectures aim to replicate the success of Large Language Models. Furthermore, state-of-the-art supervised architectures like PatchTST \citep{nie2023time} utilize patching mechanisms to significantly reduce computational complexity, theoretically allowing the processing of extended historical sequences without the quadratic cost of standard attention. However, these approaches typically operate under the assumption that maximizing lookback windows monotonically improves predictive accuracy. The proposed analysis indicates that in stochastic environments, distant history often introduces high-frequency volatility rather than actionable signal. Consequently, indiscriminate context scaling frequently leads to performance degradation due to the "noise accumulation" bottleneck, a phenomenon less prevalent in the semantic domains where these architectures originated.

\subsection{Context-Informed Forecasting and Retrieval}
Integration of auxiliary information has emerged as a critical frontier for increasing the robustness of models. \cite{williams2024context} argues that "context is key," showing that the grounding forecasts with essential textual descriptions reduces significantly the ambiguity in foundation models. Similarly, \cite{arango2025chronosx} propose ChronosX, which adapts pretrained models by consuming exogenous variables to specialize in subsequent tasks. These methods highlight the need of "context-informed" predictions. However, standard autoregressive models generally lack a mechanism for dynamically filtering out historical data according to relevance. The Retrieval-Augmented Forecasting (RAFT) framework therefore, conceptualizes historical segments not as a continuous sequence, but as a retrievable context. This approach is consistent with the philosophy of using exogenous signals but uses strict similarity-based retrieval to reject stochastic noise and effectively acts as a dynamic inductive bias.

\section{Methodology and Analysis}

\subsection{The Experimental Framework}

The experimental design rigorously evaluates the "Long Context" hypothesis on the ETTh1 (Electricity Transformer Temperature) benchmark \citep{zhou2021informer}. This dataset is ideal for stress testing  because of its high-frequency stochastic volatility, which requires the models to distinguish between semantic signals and random noise.

The study defines a fixed prediction horizon of $H=96$ steps while a varying the input lookback window $L \in \{720, 1440, 3000\}$ is used to isolate the effects of context scaling. The evaluation uses Mean Squared Error (MSE) and Mean Absolute Error (MAE) as the primary metrics. The analysis compares three different architectural paradigms to contrast \textit{Continuous Context} versus \textit{Retrieved Context}:

\begin{itemize}
    \item \textbf{Vanilla Transformer (Baseline):} The canonical encoder-decoder architecture ($d_{model}=512, n_{heads}=8$) establishes basic attention scaling properties.
    \item \textbf{PatchTST (SOTA Long-Context):} The state-of-the-art architecture \citep{nie2023time} uses channel independence and patching ($P=16, S=8$). This model represents the scaling hypothesis, which theoretically maximizes the use of context while minimizing computational complexity.
    \item \textbf{RAFT (Retrieval Baseline):} The Retrieval-Augmented Forecasting Transformer (RAFT) \citep{han2025raft} serves as the representative baseline for retrieval-based methods. Unlike the scaling models, RAFT uses a constrained lookback window ($L=720$) augmented by selectively retrieved historical segments (top-$k$ via cosine similarity) to test the hypothesis that the quality of the context outweighs the quantity.
\end{itemize}

\subsection{The Inverse Scaling Law}
Contrary to the "scaling laws" seen in Natural Language Processing, the empirical analysis reveals a different \textbf{Inverse Scaling Law} for stochastic time series. As shown in Figure \ref{fig:bar_chart}, performance does not improve with the length of the context; on the contrary, the error rate increases significantly as the window size increases.

\begin{figure}[ht]
    \centering
    \includegraphics[width=0.85\textwidth]{figures/graph3_context_vs_mse.png}
    \caption{The Inverse Scaling Law on ETTh1. While the retrieval-based RAFT (Green) maintains low error (MSE 0.379) with a fixed window, extending the context of PatchTST (Orange) to 3000 steps degrades performance by 68\% (MSE 0.647), and the Vanilla Transformer (Red) collapses completely.}
    \label{fig:bar_chart}
\end{figure}

In particular, the extension of the Vanilla Transformer context window from 720 to 3,000 steps leads to a catastrophic performance degradation and increases MSE by around \textbf{138\%}. This observation suggests that standard self-attention fails to distinguish signal from noise over long horizons.

In addition, the optimized PatchTST architecture, designed specifically for long-term forecasting, is highly sensitive to this noise accumulation. When the window extended to 3,000 steps, performance drops by \textbf{68\%} (MSE $0.385 \rightarrow 0.647$). These findings challenge the prevailing assumption that patching mechanisms alone are sufficient to filter out the noise that is inherent in large historical contexts.

% yet to be paraphared
\subsection{Generalization Across Domains}
To ensure the Inverse Scaling Law is not an artifact of the ETTh1 dataset, we validated the hypothesis on two additional benchmarks: ETTh2  \citep{zhou2021informer} (Electricity) and **Exchange Rate** (Financial).

\begin{table}[h]
\centering
\caption{Cross-Domain Validation: Degradation of Long Context}
\label{tab:cross_domain}
\small
\begin{tabular}{llcc}
\toprule
\textbf{Dataset} & \textbf{Model} & \textbf{MSE $\downarrow$} & \textbf{Degradation} \\
\midrule
\multirow{2}{*}{ETTh2} & PatchTST (720) & 0.307 & --- \\
& PatchTST (3000) & 0.533 & \textbf{+73.4\%} \\
\midrule
\multirow{2}{*}{Exchange} & PatchTST (720) & 0.093 & --- \\
& PatchTST (3000) & 0.350 & \textbf{+276.8\%} \\
\bottomrule
\end{tabular}
\end{table}

As shown in Table \ref{tab:cross_domain}, the degradation is consistent and severe. Notably, on the **Exchange Rate** dataset, which is characterized by high financial volatility, the long-context model degrades by **276\%**. This confirms that the noise accumulation bottleneck is a fundamental property of stochastic time series, intensifying in domains with lower signal-to-noise ratios.

\subsection{Mechanism Analysis: Stochastic Noise Accumulation}
The study attributes this degradation to a phenomenon termed \textbf{Stochastic Noise Accumulation}. In semantic areas such as NLP, distant tokens often retain a high degree of information density. However, in the ETTh1 dataset, distant historical points are often uncorrelated stochastic fluctuations rather than actionable signals.

As a result, the context window ($L$) increases, the attention mechanism distributes the Softmax scores to a much larger denominator. The probability mass given to the most relevant recent signals is diluted by the accumulated noise of thousands of irrelevant timesteps. The model effectively over-smooths its internal representation and loses the high-frequency conceptual focus that is necessary for accurate prediction. This observation aligns with the "Context is Key" findings of \cite{williams2024context} which suggests that more data does not mean better context when the signal-to-noise ratio is low.

\subsection{The Superiority of Selective Retrieval}
The performance of the RAFT baseline highlights the structural advantage of selective retrieval over blind scaling. By utilizing a specific \textbf{Context-Informed Inductive Bias}, RAFT filters the history before the attention mechanism processes it.

\begin{figure}[ht]
    \centering
    \includegraphics[width=0.85\textwidth]{figures/graph1_mse_comparison.png}
    \caption{Performance degradation trend. The error rate (Y-axis) increases monotonically as the Context Length (X-axis) expands from 720 to 3000, illustrating the Stochastic Noise Accumulation.}
    \label{fig:line_chart}
\end{figure}

This structural prior yields superior performance. As shown in Figure \ref{fig:line_chart}, the RAFT achieves an \textbf{MSE of 0.379} using only a 720-step window. This result surpasses the best PatchTST configuration (MSE 0.385) and significantly better than the PatchTST 3000 configuration (MSE 0.647). By treating historical segments as \textbf{dynamic exogenous variables} \citep{arango2025chronosx}, RAFT stabilizes the forecast without the penalty of continuous context scaling and demonstrates that retrieval provides a more reliable mechanism for noise rejection than window extension.

\section{Discussion}

\subsection{The Stochastic Context Dilemma}
The divergence between the NLP scaling laws and the observed "Inverse Scaling Law" for time series warrants a structural examination of the attention mechanism. In large language models, the "Needle in a Haystack" problem can be solved because semantic tokens (such as a specific entity name) retain a high signal-to-noise ratio, irrespective of their position in the sequence. In contrast, in stochastic time series benchmarks such as ETTh1, historical points far from the forecast horizon often have a close to zero correlation with the target distribution.

This analysis indicates that the extension of context window indiscriminately increases \textit{Attention Entropy}. When the attention mechanism calculates the scores of the Softmax over a 3,000-step history, the probability mass spreads thinly across thousands of irrelevant data points. This "dilution" effect drowns out the localized signals that are necessary for accurate short-term forecasting. So the architectural assumption that more history equals better context holds valid only if the historical data has a high semantic density, a condition that is rarely met in raw stochastic volatility.

\subsection{The Efficiency-Accuracy Trade-off}
Beyond accuracy, the computational costs of scaling the context scaling continuously are prohibitive. Our analysis reveals a stark disconnect between resource consumption and performance:

\begin{itemize}
    \item \textbf{Long Context (Inefficient):} The Vanilla Transformer with a 3,000-step window required 83.47 minutes to train, yet achieved the worst performance (MSE 1.323).
    \item \textbf{Retrieval (Efficient):} RAFT required only 2.13 minutes to train with an approximate 40$\times$ speedup while achieving the best performance (MSE 0.379).
\end{itemize}

This finding contradicts the "Scaling Law" ethos where increasing computing power yields increased intelligence. For time series forecasting, the processing of thousands of irrelevant historical steps is not only computationally wasteful, but also impairs the accuracy of the data.

\subsection{Retrieval as a Denoising Inductive Bias}
The better performance of the RAFT demonstrates demonstrates the effectiveness of Context-Informed Inductive Biases. By only collecting the most similar historical segments with a top-$k$ value, the model effectively implements a "hard attention" mechanism before the data is processed by the standard "soft attention" layer. This two-step process acts as a noise-rejection filter:
\begin{enumerate}
    \item \textbf{Filter:} The retrieval step discards uncorrelated history, ensuring that the input window contains only high-fidelity reference patterns.
    \item \textbf{Process:} The forecast model then treats these purified signals as \textit{dynamic exogenous variables} \citep{arango2025chronosx} rather than continuous noise.
\end{enumerate}
Consequently, this structural prior allows the model to exploit long-term historical dependencies without suffering from the noise accumulation penalty of scaling of the window.

\section{Conclusion and Future Work}

The proposed paper presents a critical assessment of the "Long Context" hypothesis in time series forecasting. A rigorous empirical analysis of the ETTh1 shows a stark \textbf{Inverse Scaling Law}: performance of modern architectures such as PatchTST declines by more than 68\% as the context length expands from 720 to 3,000 steps. These findings suggest that in stochastic domains brute force context scaling introduces a detrimental noise accumulation which outweighs the theoretical benefits of a longer history.

In addition, the evaluation of the Retrieval-Augmented Forecasting (RAFT) \citep{han2025raft} baseline shows that selective retrieval is a structurally better option. By achieving an MSE of 0.379 with a limited window, the retrieval approach outperforms the computationally costly long-context models. This confirms the hypothesis that future time series foundation models must give priority to the quality of the context over the quantity of the context.

Future research should explore the integration of this retrieval mechanism directly into the pretraining objective of foundation models. Instead of relying on static context windows, next generation architectures could use dynamic retrieval heads that learn to query historical databases for critical context, in line with the "Context is Key" philosophy \citep{williams2024context} allowing for robust zero-shot generalization in noisy environments.

\bibliography{iclr2026_conference}
\bibliographystyle{iclr2026_conference}

\appendix

\section{Experimental Configuration}

This section provides complete hyperparameter settings and training configurations to ensure reproducibility of our results.

\subsection{Model Hyperparameters}

Table~\ref{tab:hyperparams} presents the complete configurations for all evaluated models. All models were trained with identical optimization settings to ensure fair comparison.

\begin{table}[h]
\centering
\caption{Complete Hyperparameter Configurations}
\label{tab:hyperparams}
\small
\begin{tabular}{lccc}
\toprule
\textbf{Parameter} & \textbf{Vanilla TF} & \textbf{PatchTST} & \textbf{RAFT} \\
\midrule
seq\_len & 720/1440/3000 & 720/3000 & 720 \\
pred\_len & 96 & 96 & 96 \\
label\_len & 48 & 48 & 48 \\
$d_{\text{model}}$ & 512 & 128 & 128 \\
$n_{\text{heads}}$ & 8 & 16 & 8 \\
$e_{\text{layers}}$ & 2 & 3 & 2 \\
$d_{\text{layers}}$ & 1 & 1 & 1 \\
$d_{\text{ff}}$ & 2048 & 256 & 256 \\
dropout & 0.1 & 0.1 & 0.1 \\
patch\_len & --- & 16 & --- \\
stride & --- & 8 & --- \\
top-$k$ retrieval & --- & --- & 5 \\
\midrule
learning\_rate & 0.0001 & 0.0001 & 0.0001 \\
optimizer & Adam & Adam & Adam \\
batch\_size & 32 & 32 & 32 \\
epochs & 10 & 10 & 10 \\
patience & 3 & 3 & 3 \\
loss & MSE & MSE & MSE \\
\bottomrule
\end{tabular}
\end{table}

\subsection{Training Protocol}

All models were trained using the Adam optimizer with a fixed learning rate of 0.0001. We employed early stopping with a patience of 3 epochs based on validation MSE. Random seeds were fixed to 2021 for weight initialization to ensure deterministic results.

\section{Dataset Details}

We evaluate our hypothesis across three distinct time series benchmarks to demonstrate the generalizability of the inverse scaling law.

\subsection{ETTh1 Dataset}

The ETTh1 (Electricity Transformer Temperature) dataset contains hourly measurements from an electricity transformer with 14,400 timesteps spanning 600 days. The dataset includes 7 features: HUFL, HULL, MUFL, MULL, LUFL, LULL, and OT (Oil Temperature). We follow the standard 60/20/20 split for training, validation, and testing. All features are normalized using StandardScaler.

\subsection{Cross-Domain Datasets}

ETTh2 follows the same structure as ETTh1 but records measurements from a different transformer station. The Exchange Rate dataset contains daily exchange rates for 8 currencies over 7,588 days. These datasets provide validation across electricity and financial domains with varying signal-to-noise characteristics.

\section{Complete Results}

Table~\ref{tab:complete_results} presents comprehensive results including both MSE and training time metrics.

\begin{table}[h]
\centering
\caption{Complete Results on ETTh1}
\label{tab:complete_results}
\small
\begin{tabular}{lccc}
\toprule
\textbf{Model} & \textbf{Seq Len} & \textbf{MSE $\downarrow$} & \textbf{Time (min)} \\
\midrule
RAFT & 720 & \textbf{0.379} & 2.13 \\
PatchTST & 720 & 0.385 & 2.90 \\
PatchTST & 3000 & 0.647 & 18.33 \\
Vanilla TF & 720 & 0.441 & 7.42 \\
Vanilla TF & 1440 & 0.484 & 22.00 \\
Vanilla TF & 3000 & 1.323 & 83.47 \\
\bottomrule
\end{tabular}
\end{table}

\subsection{Degradation Analysis}

Table~\ref{tab:degradation} quantifies the performance degradation across different context lengths and datasets.

\begin{table}[h]
\centering
\caption{Performance Degradation Across Contexts}
\label{tab:degradation}
\small
\begin{tabular}{lccc}
\toprule
\textbf{Model} & \textbf{Dataset} & \textbf{720 $\rightarrow$ 3000} & \textbf{Degradation} \\
\midrule
Vanilla TF & ETTh1 & 0.441 $\rightarrow$ 1.323 & +200\% \\
PatchTST & ETTh1 & 0.385 $\rightarrow$ 0.647 & +68\% \\
PatchTST & ETTh2 & 0.307 $\rightarrow$ 0.533 & +73\% \\
PatchTST & Exchange & 0.093 $\rightarrow$ 0.350 & +277\% \\
\bottomrule
\end{tabular}
\end{table}

The degradation pattern remains consistent across all tested datasets. Financial data exhibits the most severe degradation, reflecting higher stochastic volatility in currency markets.

\section{Computational Environment}

All experiments were conducted on cloud GPU infrastructure with 32GB RAM and 16 CPU cores. We used Python 3.9, PyTorch 2.0.1, and CUDA 11.8. The complete software stack includes numpy 1.24, pandas 2.0, and scikit-learn 1.3.

Training times scale quadratically with context length. The Vanilla Transformer with 3000-step context required 83.47 minutes per run, while RAFT completed training in 2.13 minutes, achieving a 40$\times$ speedup with superior accuracy.

\section{Reproducibility}

Our experimental setup ensures full reproducibility. All datasets are publicly available from the Informer repository. We use standard train-validation-test splits of 60/20/20 and apply per-feature StandardScaler normalization. Random seeds are fixed to 2021 for all experiments. Early stopping monitors validation MSE with patience of 3 epochs.

To reproduce our results, download the datasets, apply the configurations in Table~\ref{tab:hyperparams}, and train with the specified optimizer settings. All models use standard PyTorch implementations without custom modifications.

\section{Limitations}

Our analysis focuses on univariate forecasting with a fixed 96-step horizon. Results may differ for longer prediction horizons or multivariate tasks. We evaluate on regular time series; irregular sampling patterns may exhibit different scaling behaviors.

The retrieval mechanism in RAFT relies on cosine similarity, which may not be optimal for all domains. Future work should explore domain-adaptive retrieval metrics and investigate hybrid architectures combining selective retrieval with moderate context windows.


\end{document}
